\bookchapter{Capítulo 33: Perspectivas futuras — El próximo recorrido del móvil}

\chapterintro{Revisa los permisos de las apps}

\booksection{: Ajustes > Privacidad > Permisos. Revisa qué permisos tienen las apps; niega los innecesarios. Una app de linterna no necesita acceso a tus contactos.}

No conectes el teléfono a Wi-Fi público desconocido

\begin{itemize}
  \item \textbf{: Los Wi-Fi públicos pueden ser puntos de interceptación de datos. Si tienes que conectarte, no hagas operaciones bancarias ni introduzcas contraseñas; usa tu red móvil en su lugar.}No instales apps de fuentes desconocidas
  \item \textbf{: activa "Buscar mi móvil" y establece una contraseña de bloqueo de pantalla robusta.}: Solo instala apps desde la App Store (iPhone) o tiendas oficiales de Android (Google Play, Huawei AppGallery, etc.). Los APK de fuentes desconocidas pueden contener malware.
  \item \textbf{No hagas clic en enlaces sospechosos}: Si recibes un SMS o un mensaje de WeChat con un enlace que no esperabas, no lo abras; puede ser un enlace de phishing que roba tus contraseñas.
\end{itemize}

\booksection{Activa el pago biométrico}

: Activa el pago con Face ID o huella en Alipay/WeChat Pay. Es más seguro que la contraseña y evita que alguien vea tu contraseña por encima del hombro.

\begin{itemize}
  \item \textbf{Establece una "contraseña familiar"}: Acuerda con tus familiares una contraseña que solo vosotros conozcáis. Si alguien llama pidiendo dinero urgente diciendo ser un familiar, pídele la contraseña; si no la sabe, es una estafa.
  \item \textbf{Mantén el sistema actualizado}: Las actualizaciones del sistema incluyen parches de seguridad. No ignores las notificaciones de actualización; instálalas lo antes posible.
  \item \textbf{IV. Diez errores comunes que debes evitar}Dejar el GPS siempre activado
\end{itemize}

\booksection{: El GPS consume batería y revela tu ubicación. Desactívalo cuando no lo necesites; actívalo solo para navegar o usar apps que lo requieran.}

No hacer copia de seguridad del teléfono

\begin{itemize}
  \item \textbf{Príorizar tiendas oficiales}: Si pierdes o rompes el teléfono, pierdes todas tus fotos, contactos y mensajes. Activa la copia de seguridad automática (iCloud / Google One / copia local del fabricante).
  \item \textbf{No des permisos a la ligera}Usar el teléfono mientras se carga
  \item \textbf{Apéndice B: Referencia de ejercicios prácticos}: Usar el teléfono mientras se carga genera calor extra y acelera el envejecimiento de la batería. Si tienes que usarlo, al menos no juegues ni grabes vídeos.
\end{itemize}

\booksection{Dejar que la batería llegue a 0\%}

Ejercicio de software

\begin{itemize}
  \item \textbf{: Las baterías de litio no les gusta llegar a 0\% ni quedarse al 100\% mucho tiempo. Intenta cargar el teléfono cuando llegue al 20-30\% y desenchufarlo al 80-90\%.}Dejar el brillo de la pantalla al máximo
  \item \textbf{: El brillo máximo consume batería y daña los ojos. Activa el brillo automático o ajústalo manualmente al nivel mínimo cómodo.}No limpiar el objetivo de la cámara
  \item \textbf{: El objetivo de la cámara se ensucia con el uso y las fotos salen borrosas. Límpialo con un paño suave antes de tomar fotos importantes.}Instalar demasiadas apps "por si acaso"
  \item : configura la limpieza de almacenamiento de WeChat; ruta habitual: [Yo] -> [Configuracion] -> [General] -> [Almacenamiento].
  \item Ejercicio Intermedio
\end{itemize}

\booksection{: Cada app instalada consume almacenamiento, memoria y batería. Desinstala las apps que no uses; puedes reinstalarlas cuando las necesites.}

No leer los permisos al instalar una app

\begin{enumerate}
  \item Truco universal
  \item : ante cualquier problema, reinicia primero! Esta es la llave maestra que resuelve más de la mitad de los pequenos problemas del móvil.
  \item Apndice C: Tabla de comparacion de terminologia "jerga técnica"
  \item Termino que puedes escuchar
  \item En realidad se refiere a...
\end{enumerate}

\begin{tcbpractice}{Explicacion en una frase}
: Mucha gente aprueba todos los permisos sin pensar. Lee los permisos que pide la app; si una app pide permisos que no necesita, no la instales.

\begin{itemize}
  \item El cerebro y corazon del móvil
  \item Componente clave que determina la rapidez del móvil y si los juegos van fluidos.
  \item RAM (memoria)
\end{itemize}

\end{tcbpractice}

\begin{tcbsummary}{Resumen del capítulo}
\begin{itemize}
  \item Dejar el teléfono sin código de desbloqueo
  \item : Si pierdes el teléfono y no tiene código de desbloqueo, cualquiera puede acceder a tus fotos, mensajes y cuentas. Pon siempre un código de desbloqueo, Face ID o huella.
  \item Compartir la ubicación en tiempo real en redes sociales
  \item : Publicar fotos con ubicación en tiempo real en WeChat Moments o Xiaohongshu permite a los estafadores saber dónde estás. Comparte la ubicación después de irte, no en tiempo real.
  \item V. Diez pequeños trucos que mejoran la vida
\end{itemize}

\end{tcbsummary}
